% Relatório de estágio supervisionado em Biblioteconomia
% Modelo do ÉguaTex (https://eguatex.github.io), em domínio público pela licença CC0 1.0:
% use, modifique e distribua sem pedir permissão e sem citar a origem.
%
% Compile com pdfLaTeX e BibTeX (o ÉguaTex faz isso sozinho).
% Arquivos: main.tex (este texto) e referencias.bib (bibliografia).

\documentclass[12pt,a4paper]{report}

% ---------- idioma, fontes e página ----------
\usepackage[T1]{fontenc}
\usepackage[utf8]{inputenc}
\usepackage[brazilian]{babel}
\usepackage{tgtermes}
\usepackage{microtype}
\usepackage[a4paper,left=3cm,top=3cm,right=2cm,bottom=2cm]{geometry}
\usepackage{setspace}
\onehalfspacing
\setlength{\parindent}{1.25cm}

% ---------- tabelas, quadros e cores ----------
\usepackage{amssymb}
\usepackage[table]{xcolor}
\usepackage{array}
\usepackage{tabularx}
\usepackage{booktabs}
\usepackage{caption}
\captionsetup{font=small, labelfont=bf, justification=centering}
% A ABNT chama de quadro a tabela de texto, com bordas; troque por Tabela se preferir.
\captionsetup[table]{name=Quadro}
\usepackage{enumitem}
\setlist{nosep, leftmargin=1.25cm}

\definecolor{marcado}{RGB}{46,110,82}   % semanas com atividade no cronograma
\definecolor{fundo}{RGB}{232,241,236}   % cabeçalho dos quadros

% ---------- títulos das seções: CAIXA ALTA, como pede a ABNT ----------
\renewcommand{\thesection}{\arabic{section}}
\usepackage{titlesec}
\titleformat{\section}{\normalsize\bfseries}{\thesection}{0.6em}{\MakeUppercase}
\titleformat{\subsection}{\normalsize\bfseries}{\thesubsection}{0.6em}{}
\titlespacing*{\section}{0pt}{1.5em}{0.8em}

% ---------- links e referências no padrão ABNT ----------
% O hyperref vem antes do abntex2cite, senão as citações quebram.
\usepackage[hidelinks]{hyperref}
\usepackage[alf]{abntex2cite}
\addto\captionsbrazilian{\renewcommand{\bibname}{Referências}\renewcommand{\contentsname}{Sumário}}

% ---------- instruções de preenchimento ----------
% Os textos em cinza explicam o que escrever em cada parte.
% Troque \instrucoestrue por \instrucoesfalse para esconder todos de uma vez.
\newif\ifinstrucoes
\instrucoestrue
\newcommand{\instrucao}[1]{\ifinstrucoes{\color{gray}\small\itshape #1\par}\fi}

% ---------- dados do estágio: preencha aqui ----------
\newcommand{\universidade}{Universidade Federal Exemplo}
\newcommand{\unidade}{Faculdade de Biblioteconomia}
\newcommand{\curso}{Curso de Biblioteconomia}
\newcommand{\estagiario}{Nome do autor}
\newcommand{\titulo}{Relatório de estágio supervisionado}
\newcommand{\concedente}{Biblioteca Central da Universidade Federal Exemplo}
\newcommand{\supervisor}{Nome do supervisor ou da supervisora (CRB 0/0000)}
\newcommand{\orientador}{Nome do orientador ou da orientadora}
\newcommand{\periodo}{dd/mm/aaaa a dd/mm/aaaa}
\newcommand{\cargahoraria}{180 horas}
\newcommand{\local}{Cidade}
\newcommand{\data}{Ano}

\hypersetup{pdftitle={\titulo}, pdfauthor={\estagiario}}

% ---------- cronograma: \m marca uma semana com atividade ----------
\newcommand{\m}{\cellcolor{marcado}}
\newcolumntype{S}{>{\centering\arraybackslash}p{0.42cm}}

\begin{document}

% ======================================================================
% Capa
% ======================================================================
\begin{titlepage}
  \centering
  {\bfseries\MakeUppercase{\universidade}\par}
  {\MakeUppercase{\unidade}\par}
  {\MakeUppercase{\curso}\par}
  \vspace{3cm}
  {\bfseries\MakeUppercase{\estagiario}\par}
  \vfill
  {\large\bfseries\MakeUppercase{\titulo}\par}
  \vspace{0.6em}
  {\concedente\par}
  \vfill
  \hfill
  \begin{minipage}{0.55\textwidth}
    \singlespacing\small
    Relatório apresentado à disciplina Estágio Supervisionado do \curso{} da \universidade{}, como requisito parcial para a obtenção do grau de Bacharel em Biblioteconomia.\par
    \medskip
    Orientação: \orientador.
  \end{minipage}
  \vfill
  {\bfseries\MakeUppercase{\local}\par\data\par}
\end{titlepage}

\tableofcontents
\clearpage

% ======================================================================
\section{Identificação do estágio}
% ======================================================================
\instrucao{Preencha o quadro com os dados do termo de compromisso. Ele resume quem, onde, quando e com quantas horas o estágio foi feito.}

\begin{table}[htbp]
  \centering
  \caption{Dados do estágio.}
  \label{tab:identificacao}
  \renewcommand{\arraystretch}{1.3}
  \begin{tabularx}{\textwidth}{>{\bfseries}p{4.2cm} X}
    \toprule
    \rowcolor{fundo} Item & Descrição \\
    \midrule
    Estagiário(a)         & \estagiario \\
    Unidade concedente    & \concedente \\
    Setor                 & Processamento técnico e serviço de referência \\
    Supervisão no local   & \supervisor \\
    Orientação na UFE     & \orientador \\
    Período               & \periodo \\
    Carga horária         & \cargahoraria{}, em quatro horas diárias \\
    \bottomrule
  \end{tabularx}
\end{table}

% ======================================================================
\section{Introdução}
% ======================================================================
\instrucao{Em dois ou três parágrafos: o que é o estágio, por que ele faz parte da formação e quais objetivos você levou para a unidade. Exemplo: aplicar a catalogação descritiva e conhecer a rotina do atendimento ao usuário.}

O estágio supervisionado aproxima a formação acadêmica da rotina das unidades de informação. Neste relatório, descrevem-se as atividades realizadas na \concedente{}, organizadas segundo as cinco leis da Biblioteconomia, que orientam o serviço para o uso do acervo e para o atendimento ao leitor \cite{ranganathan2009}.

Os objetivos do estágio foram conhecer os fluxos de processamento técnico, praticar a catalogação e a classificação de materiais e acompanhar o serviço de referência. O texto segue a NBR 14724, que trata da apresentação de trabalhos acadêmicos \cite{nbr14724}.

% ======================================================================
\section{Caracterização da unidade concedente}
% ======================================================================
\instrucao{Descreva a biblioteca em meia página: vínculo institucional, público atendido, tamanho do acervo, equipe, serviços oferecidos e sistema de gestão usado. Use dados fornecidos pela própria unidade.}

A biblioteca atende estudantes, docentes e servidores da universidade e a comunidade externa para consulta local. Descreva aqui o acervo, os setores e a equipe, com números informados pela direção da unidade.

% ======================================================================
\section{Atividades desenvolvidas}
% ======================================================================
\instrucao{Liste as atividades por setor e mostre no cronograma em que semanas cada uma ocorreu. Depois comente as mais importantes, com exemplos do que você fez.}

\subsection{Cronograma de atividades}

O Quadro~\ref{tab:cronograma} mostra a distribuição das atividades ao longo das doze semanas do estágio. Cada célula preenchida indica uma semana em que a atividade foi realizada.

% Para marcar uma semana, escreva \m na célula correspondente.
\begin{table}[htbp]
  \centering
  \caption{Cronograma das atividades do estágio, por semana.}
  \label{tab:cronograma}
  \small
  \setlength{\tabcolsep}{2.5pt}
  \renewcommand{\arraystretch}{1.35}
  \begin{tabular}{|p{5.6cm}|*{12}{S|}}
    \hline
    \rowcolor{fundo}
    \textbf{Atividade} & 1 & 2 & 3 & 4 & 5 & 6 & 7 & 8 & 9 & 10 & 11 & 12 \\
    \hline
    Reconhecimento dos setores              & \m & \m & & & & & & & & & & \\
    \hline
    Registro e preparo físico do material   & & \m & \m & \m & & & & & & & & \\
    \hline
    Catalogação descritiva                  & & & \m & \m & \m & \m & \m & & & & & \\
    \hline
    Classificação e indexação de assuntos   & & & & & \m & \m & \m & \m & & & & \\
    \hline
    Atendimento no serviço de referência    & & & & & & & & \m & \m & \m & \m & \\
    \hline
    Ação de divulgação do acervo            & & & & & & & & & \m & \m & & \\
    \hline
    Redação do relatório                    & & & & & & & & & & \m & \m & \m \\
    \hline
  \end{tabular}
\end{table}

\subsection{Processamento técnico}
\instrucao{Conte o que foi feito e com que instrumentos: código de catalogação, tabela de classificação, vocabulário controlado e sistema. Dê números, como a quantidade de registros criados.}

Descreva as etapas de registro, catalogação e classificação, com a quantidade aproximada de itens tratados e os instrumentos usados.

\subsection{Serviço de referência}
\instrucao{Relate os tipos de pergunta recebidos e como foram respondidos. Exemplo: orientação para normalizar referências segundo a NBR 6023.}

Relate o atendimento ao usuário e as orientações de normalização, que seguem a NBR 6023 para referências \cite{nbr6023} e a NBR 10520 para citações \cite{nbr10520}.

% A avaliação fica em página própria, para o supervisor imprimir e assinar.
\clearpage
% ======================================================================
\section{Avaliação do supervisor}
% ======================================================================
\instrucao{Este quadro é preenchido e assinado pelo supervisor da unidade. Marque um X na coluna que corresponde ao desempenho em cada critério (1 = insuficiente, 5 = excelente).}

\begin{center}
  \captionsetup{hypcap=false}% quadro fora de float, sem âncora de link
  \captionof{table}{Avaliação do desempenho do estagiário pelo supervisor.}
  \label{tab:avaliacao}
  \renewcommand{\arraystretch}{1.4}
  \begin{tabular}{|p{7.6cm}|*{5}{>{\centering\arraybackslash}p{0.9cm}|}}
    \hline
    \rowcolor{fundo}
    \textbf{Critério} & \textbf{1} & \textbf{2} & \textbf{3} & \textbf{4} & \textbf{5} \\
    \hline
    Assiduidade e pontualidade              & $\square$ & $\square$ & $\square$ & $\square$ & $\square$ \\ \hline
    Domínio das técnicas de processamento   & $\square$ & $\square$ & $\square$ & $\square$ & $\square$ \\ \hline
    Qualidade do atendimento ao usuário     & $\square$ & $\square$ & $\square$ & $\square$ & $\square$ \\ \hline
    Iniciativa e autonomia                  & $\square$ & $\square$ & $\square$ & $\square$ & $\square$ \\ \hline
    Relacionamento com a equipe             & $\square$ & $\square$ & $\square$ & $\square$ & $\square$ \\ \hline
    Cumprimento das normas da unidade       & $\square$ & $\square$ & $\square$ & $\square$ & $\square$ \\ \hline
  \end{tabular}
\end{center}

\noindent Comentários do supervisor:\par
\vspace{0.3em}
\noindent\rule{\textwidth}{0.4pt}\par\vspace{0.9em}
\noindent\rule{\textwidth}{0.4pt}\par\vspace{0.9em}
\noindent\rule{\textwidth}{0.4pt}\par

\vspace{3em}
\begin{center}
  \begin{tabular}{c@{\hspace{2cm}}c}
    \rule{6cm}{0.4pt} & \rule{5cm}{0.4pt} \\
    \small Assinatura do supervisor & \small Data
  \end{tabular}
\end{center}

% ======================================================================
\section{Considerações finais}
% ======================================================================
\instrucao{Em um ou dois parágrafos: o que o estágio acrescentou à sua formação, as dificuldades encontradas e uma sugestão concreta para a unidade ou para o curso.}

O estágio permitiu aplicar os conteúdos de representação descritiva e temática em um acervo real. Registre aqui as aprendizagens, as dificuldades e as sugestões de melhoria.

% ======================================================================
% Referências: seção sem número, na mesma página, listada no sumário.
% ======================================================================
\begingroup
\let\chapter\section
\phantomsection
\addcontentsline{toc}{section}{\bibname}
\bibliography{referencias}
\endgroup

\end{document}
